\documentclass[11pt,a4paper]{article}

\usepackage[margin=2.2cm]{geometry}
\usepackage{booktabs}
\usepackage{amsmath}
\usepackage{listings}
\usepackage{xcolor}
\usepackage{hyperref}
\usepackage{longtable}
\usepackage{fancyhdr}
\usepackage{enumitem}

\definecolor{codebg}{RGB}{246,247,249}
\definecolor{codekw}{RGB}{20,101,95}
\definecolor{annot}{RGB}{90,90,90}

\lstdefinelanguage{GraphEasy}{
  keywords={
    graph, graphs, motifs, directed, nodes, edges, file, weights, TRUE,
    for, each, vertex, edge, neighbor, of, in, where, motif, add, remove,
    to, from, draw, show, query, print, set, fn, int, real, bool, string,
    void, while, if, else, return, swap, sleep, union, intersect, contains,
    layout, vertices, labels, color, size, categorical, continuous,
    degree, connected, with, cycle, timer, INF, numVertices, numEdges,
    numMotifs, numGraphs, hasEdge, weight, out, in
  },
  keywordstyle=\color{codekw}\bfseries,
  sensitive=false,
  morecomment=[l]{//},
  morecomment=[s]{/*}{*/},
  morestring=[b]",
}

\lstset{
  language=GraphEasy,
  basicstyle=\ttfamily\small,
  backgroundcolor=\color{codebg},
  frame=single,
  framerule=0pt,
  breaklines=true,
  showstringspaces=false,
  tabsize=2,
  xleftmargin=8pt,
  xrightmargin=8pt,
  aboveskip=8pt,
  belowskip=8pt,
}

\newcommand{\ann}[1]{\par\vspace{2pt}{\small\textcolor{annot}{\textit{#1}}}\vspace{4pt}}
\newcommand{\kw}[1]{\texttt{\textcolor{codekw}{#1}}}
\newcommand{\sym}[1]{\texttt{#1}}

\hypersetup{colorlinks=true, linkcolor=codekw, urlcolor=codekw}

\pagestyle{fancy}
\fancyhf{}
\rhead{\small GraphEasy graph grammar}
\lhead{\small Language reference}
\cfoot{\thepage}

\title{\textbf{GraphEasy Graph Grammar}\\[0.3em]
\large Declarations, traversal, comprehension, motifs, and drawing}
\author{Graph-language project \quad (\texttt{Base.g4})}
\date{2026-08-27}

\begin{document}
\maketitle

\begin{abstract}
\noindent
This note documents the \emph{graph-related} surface syntax of GraphEasy:
how to declare graphs, load edge lists, traverse vertices and edges,
filter subgraphs, search motifs, mutate structure, query, print, and draw.
Each construct is tied to the ANTLR rule in \texttt{Base.g4} and illustrated
with a short \texttt{.graph} example. Non-graph statements (scalar arithmetic,
arrays, generic control flow) are omitted except where they interact with
graphs.
\end{abstract}

\tableofcontents
\newpage

%==============================================================================
\section{Programs and \texttt{.graph} files}
%==============================================================================

A GraphEasy program is a sequence of statements and optional function
definitions, terminated by end-of-file:
\[
  \texttt{program} \;:\; (\texttt{statement} \mid \texttt{function})^* \;\texttt{EOF}
\]
Programs live in files with extension \texttt{.graph}. A typical file
declares one or more graphs, runs an algorithm over them, then prints or
draws results.

%==============================================================================
\section{Graph declaration (\texttt{graphDef})}
%==============================================================================

The central construct is \kw{graph}. Grammar rule:
\[
  \texttt{graphDef} \;:\;
  \texttt{graph}\;\mathit{graphID}\;\{\;\mathit{graphProperty}^*\;
  \mathit{nodes}?\;\mathit{graphProperty}^*\;\mathit{edges}?\;
  \mathit{graphProperty}^*\;\texttt{TRUE}?\;\mathit{graphProperty}^*\;\}\;\texttt{;}
\]

\subsection{Minimal inline graph}

\begin{lstlisting}
graph G {
  nodes: 0, 1, 2, 3;
  edges: 0->1, 1->2, 2->3;
};
\end{lstlisting}

\ann{%
  \textbf{\kw{graph G}} --- declares a graph named \texttt{G}.\\
  \textbf{\kw{nodes:}} --- optional explicit vertex list (integers).\\
  \textbf{\kw{edges:}} --- comma-separated \sym{u->v} pairs.\\
  \textbf{\sym{;}} after \texttt{\}} --- statement terminator (required).
}

Vertices may be omitted when they appear only in edges; the runtime
infers node ids from edge endpoints.

\subsection{Loading edges from a file}

\begin{lstlisting}
graph G {
  edges: file "edgelist.txt";
};
\end{lstlisting}

\ann{%
  \textbf{\kw{file "..."}} --- replaces inline \kw{edges:} list.\\
  Unweighted file: two columns \texttt{u v} per line.\\
  Weighted file: three columns \texttt{u v w} (see weighted mode below).
}

\subsection{Directed graphs}

Graphs are \emph{undirected by default}. An input line \texttt{u v} creates
both \sym{u->v} and \sym{v->u} unless direction is requested.

\begin{lstlisting}
graph G {
  directed: true;
  edges: file "directed_edgelist.txt";
};
\end{lstlisting}

\ann{%
  \textbf{\kw{directed: true;}} --- graph property (\texttt{graphProperty}).\\
  Each file line \texttt{u v} becomes a single directed edge \sym{u->v}.\\
  Motif arrows \sym{->} and neighbor direction (\kw{out neighbor},
  \kw{in neighbor}) are meaningful only on directed graphs.
}

\subsection{Weighted graphs}

Two equivalent forms mark a graph as weighted:

\begin{lstlisting}
graph G {
  edges: file "weighted.txt";
  TRUE
};

graph W {
  weights: true;
  directed: true;
  edges: file "weighted.txt";
};
\end{lstlisting}

\ann{%
  \textbf{\kw{TRUE}} (legacy) or \textbf{\kw{weights: true;}} --- read third
  column from edge file.\\
  \textbf{\kw{weight(G, u, v)}} returns that weight at runtime; unweighted
  edges behave as weight~1.\\
  Extended modes: \kw{weights:positive;}, \kw{weights:zero;},
  \kw{weights:negative;} assign a fixed sign to every loaded edge
  (used in signed-motif biology-style networks).
}

\subsection{Grammar mapping}

\begin{center}
\begin{tabular}{@{}ll@{}}
\toprule
Surface syntax & Rule in \texttt{Base.g4} \\
\midrule
\texttt{graph G \{ \ldots \};} & \texttt{graphDef} \\
\texttt{directed: true;} & \texttt{graphProperty} \\
\texttt{nodes: 0, 1, 2;} & \texttt{nodes} $\rightarrow$ \texttt{nodeList} \\
\texttt{edges: 0->1, 1->2;} & \texttt{edges} $\rightarrow$ \texttt{edgeList} \\
\texttt{edges: file "g.txt";} & \texttt{edges} $\rightarrow$ \texttt{fileEdgeList} \\
\texttt{0->1} & \texttt{edge} \\
\bottomrule
\end{tabular}
\end{center}

%==============================================================================
\section{Graph traversal loops}
%==============================================================================

Graph-aware \kw{for each} loops avoid manual CSR indexing. Targets
(\texttt{loopTarget}):

\begin{center}
\begin{tabular}{@{}p{5.8cm}p{7.5cm}@{}}
\toprule
Syntax & Visits \\
\midrule
\texttt{for each vertex v in G} & every vertex id \\
\texttt{for each edge u, v in G} & every edge (ordered pair) \\
\texttt{for each neighbor v of u in G} & adjacency of \texttt{u} \\
\texttt{for each out neighbor v of u in G} & outgoing edges (directed) \\
\texttt{for each in neighbor v of u in G} & incoming edges (directed) \\
\texttt{for each graph H in S} & each graph in a \kw{graphs} collection \\
\texttt{for each motif (a,b,c) in M} & each motif role binding \\
\bottomrule
\end{tabular}
\end{center}

\subsection{Example: BFS over neighbors}

\begin{lstlisting}
graph G {
  edges: file "graph.txt";
};

int n = numVertices(G);
int parent[n];
// ... initialize parent[], frontier arrays ...

while (frontier_size > 0) {
  int v = frontier[i];
  for each neighbor u of v in G {
    if (parent[u] == -1) {
      parent[u] = v;
      // enqueue u into next frontier
    }
  }
}
\end{lstlisting}

\ann{%
  \textbf{\kw{for each neighbor u of v in G}} --- \texttt{ForEachAdj} in the
  grammar. On undirected graphs, \kw{neighbor}, \kw{out neighbor}, and
  \kw{in neighbor} coincide. On directed graphs, \kw{neighbor} follows
  outgoing edges.
}

%==============================================================================
\section{Graph sets and membership}
%==============================================================================

A graph exposes its vertex and edge sets:

\begin{lstlisting}
set V = G.nodes;      // same as G.vertices
set E = G.edges;

V.add(4);
V.remove(0);

if (2 in G) { print "vertex 2 exists"; }
if (0->1 in G) { print "edge exists"; }
\end{lstlisting}

\ann{%
  \textbf{\kw{G.nodes}} / \textbf{\kw{G.vertices}} --- vertex set
  (\texttt{GraphNodesSet}).\\
  \textbf{\kw{G.edges}} --- edge set (\texttt{GraphEdgesSet}).\\
  \textbf{\texttt{nodeID in graphID}} / \textbf{\texttt{edge in graphID}} ---
  membership tests in conditions (\texttt{NodeCheck}, \texttt{EdgeCheck}).
}

Set algebra: \kw{union}, \kw{intersect}, literals \texttt{\{0,1,2\}}.

%==============================================================================
\section{Graph comprehension}
%==============================================================================

Rule \texttt{graphComprehension}:
\[
  \mathit{ID} \;=\; \texttt{[}\;\mathit{graphID}\;
  (\texttt{where}\;\mathit{graphCondition})?\;\texttt{]}\;\texttt{;}
\]
When the left-hand type is \kw{graph}, the result is a new graph value.

\subsection{Vertex and edge filters}

\begin{lstlisting}
set P = {0, 1, 2};
graph H = [G where vertex in P];
graph Incident = [G where edge has v];
\end{lstlisting}

\ann{%
  \textbf{\kw{vertex in P}} --- induced subgraph on vertices in set \texttt{P}.\\
  \textbf{\kw{edge has v}} --- subgraph of edges incident to vertex \texttt{v}
  (typical neighbor extraction).
}

\subsection{Structural conditions}

\begin{lstlisting}
graph High = [G where degree > 2];
graph Comp = [G where connected with 0];
graph Cyclic = [G where cycle];
graph Mixed = [G where degree > 2 && connected with 0];
\end{lstlisting}

\ann{%
  \textbf{\kw{degree}} \sym{op} \texttt{INT} --- filter by vertex degree.\\
  \textbf{\kw{connected with}} \texttt{nodeID} --- keep vertices in the same
  connected component as the given vertex.\\
  \textbf{\kw{cycle}} --- retain cycle-related structure (runtime-defined).\\
  Conditions combine with \texttt{\&\&} and \texttt{||} (\texttt{GraphLogicalAnd/Or}).
}

\subsection{Motif condition inside comprehension}

A comprehension may embed a motif pattern; the result is one \emph{merged}
subgraph containing all edges that participate in \emph{any} match:

\begin{lstlisting}
graph FeedForward = [G where motif {
  a->b;
  a->c;
  b->c;
}];
\end{lstlisting}

\ann{%
  Same motif syntax as \kw{motifs}/\kw{graphs} (below), but the answer shape
  is a single \kw{graph} --- the union of motif edges --- not a list of
  matches.
}

%==============================================================================
\section{Motif search: \kw{motifs}, \kw{graphs}, and \kw{graph}}
%==============================================================================

Motifs are fixed compile-time subgraph patterns. Each edge line uses role
names (identifiers), not numeric vertex ids.

\subsection{Motif edge syntax}

\begin{center}
\begin{tabular}{@{}ll@{}}
\toprule
Motif edge & Meaning \\
\midrule
\texttt{a->b;} & positive / ordinary directed edge between roles \texttt{a}, \texttt{b} \\
\texttt{a-|b;} & negative signed edge (matches negative-weight edges in signed graphs) \\
\bottomrule
\end{tabular}
\end{center}

Grammar: \texttt{motifEdge} $\rightarrow$ \texttt{ID ('->' | '-|') ID ';'}.

\subsection{Three result shapes}

\begin{center}
\begin{tabular}{@{}p{3.2cm}p{5.5cm}p{4.5cm}@{}}
\toprule
Keyword & Example & Use \\
\midrule
\texttt{motifs} &
\texttt{motifs M = [G where motif \{...\}];} &
Count or iterate role assignments \\
\texttt{graphs} &
\texttt{graphs S = [G where motif \{...\}];} &
One extracted subgraph \emph{per} match \\
\texttt{graph} &
\texttt{graph H = [G where motif \{...\}];} &
One merged subgraph (comprehension) \\
\bottomrule
\end{tabular}
\end{center}

\subsection{Example: feed-forward loop (FFL)}

\begin{lstlisting}
graph G {
  directed: true;
  nodes: 0, 1, 2, 3, 4, 5, 6, 7, 8, 9, 10, 11, 12, 13;
  edges:
    0->1, 0->2, 1->2,
    3->4, 4->5,
    6->8, 6->9, 7->8, 7->9,
    10->11, 10->12, 11->13, 12->13;
};

motifs FFLMatches = [G where motif {
  a->b;
  a->c;
  b->c;
}];

graphs FFLGraphs = [G where motif {
  a->b;
  a->c;
  b->c;
}];

print numMotifs(FFLMatches);
print numGraphs(FFLGraphs);

for each motif (a, b, c) in FFLMatches {
  print a;
  print b;
  print c;
}

for each graph H in FFLGraphs {
  print numVertices(H);
  print numEdges(H);
}
\end{lstlisting}

\ann{%
  \textbf{Role names} \texttt{a,b,c} are arbitrary identifiers bound at each
  match.\\
  \textbf{\kw{motifs}} returns a collection of tuples; \textbf{\kw{graphs}}
  returns separate small graphs so shared edges stay duplicated per match.\\
  \textbf{\kw{numMotifs(M)}} / \textbf{\kw{numGraphs(S)}} --- built-in counters.
}

\subsection{Signed motif (incoherent FFL)}

\begin{lstlisting}
graph G {
  directed: true;
  weights: true;
  edges: file "signed_edges.txt";
};

graph IncoherentFFL = [G where motif {
  a->b;
  b-|c;
  a->c;
}];
\end{lstlisting}

\ann{%
  \sym{->} matches positive-weight edges; \sym{-|} matches negative-weight
  edges. Useful for regulatory networks where sign matters.
}

\subsection{Other named patterns (Milo-style)}

\begin{lstlisting}
graph ThreeChain = [G where motif { a->b; b->c; }];

graph BiFan = [G where motif {
  a->c; a->d; b->c; b->d;
}];

graph BiParallel = [G where motif {
  a->b; a->c; b->d; c->d;
}];
\end{lstlisting}

%==============================================================================
\section{Graph mutation}
%==============================================================================

Rules \texttt{addOperation} / \texttt{removeOperation}:

\begin{lstlisting}
add 4 to G;
add 5, 6, 7 to G;
add 2->4 to G;
add 5->6, 6->7 to G;

remove 4 from G;
remove 2->4 from G;
\end{lstlisting}

\ann{%
  \textbf{\kw{add}} \textit{targets} \textbf{\kw{to}} \textit{graph} --- nodes
  (expressions) or edges \sym{u->v}.\\
  \textbf{\kw{remove}} \textit{targets} \textbf{\kw{from}} \textit{graph} ---
  symmetric removal.
}

%==============================================================================
\section{Built-in graph functions and queries}
%==============================================================================

\subsection{Functions (selected)}

\begin{center}
\begin{tabular}{@{}llp{6.5cm}@{}}
\toprule
Call & Returns & Meaning \\
\midrule
\texttt{numVertices(G)} & \texttt{int} & $|V|$ \\
\texttt{numEdges(G)} & \texttt{int} & edge count (directed: not halved) \\
\texttt{degree(G, v)} & \texttt{int} & undirected degree, or in+out on directed \\
\texttt{outDegree(G, v)} & \texttt{int} & outgoing count \\
\texttt{inDegree(G, v)} & \texttt{int} & incoming count \\
\texttt{hasEdge(G, u, v)} & \texttt{bool} & direction-sensitive \\
\texttt{weight(G, u, v)} & \texttt{int} & edge weight (1 if unweighted) \\
\texttt{neighbors(G, v)} & \texttt{set} & neighbor set \\
\texttt{subgraph(G, S)} & \texttt{graph} & induced subgraph \\
\bottomrule
\end{tabular}
\end{center}

\subsection{Query statement}

\begin{lstlisting}
query bfs_result : "bfs" of G;
print bfs_result;
\end{lstlisting}

\ann{%
  Grammar: \texttt{query ID ':' STRING INT? 'of' graphID ';'}.\\
  The string names a built-in or runtime query; benchmark code usually
  implements algorithms explicitly instead of relying on queries.
}

%==============================================================================
\section{Printing and showing graphs}
%==============================================================================

\begin{lstlisting}
print edges of G;
print nodes of G;
print graph of G;

show G;
\end{lstlisting}

\ann{%
  \textbf{\kw{print ... of G}} --- \texttt{printgraph} (edge / node / full
  summary).\\
  \textbf{\kw{show G}} --- interactive or console display
  (\texttt{showgraph}), depending on runtime.
}

%==============================================================================
\section{Drawing graphs (\texttt{drawgraph})}
%==============================================================================

Static export to SVG, PNG, PDF, or DOT (format from filename):

\begin{lstlisting}
draw G to "network.svg";
\end{lstlisting}

Optional block \texttt{\{ drawOption* \}} controls layout and styling:

\begin{lstlisting}
draw G to "colored.png" {
  layout: "force";
  vertices {
    labels: true;
    color: categorical(classes);
    size: continuous(sizes);
  }
  edges {
    labels: true;
  }
};
\end{lstlisting}

\subsection{Layout names}

\begin{center}
\begin{tabular}{@{}ll@{}}
\toprule
\texttt{layout: "..."} & Typical use \\
\midrule
\texttt{"auto"} & default \\
\texttt{"hierarchical"} & trees, DAGs, layered flow \\
\texttt{"force"} & general networks \\
\texttt{"radial"} & hub-centric views \\
\texttt{"circular"} & small / cyclic graphs \\
\texttt{"clustered"} & grouped structure \\
\bottomrule
\end{tabular}
\end{center}

\subsection{Vertex and edge draw options}

\begin{itemize}[nosep]
  \item \kw{labels: true|false;} --- show vertex ids or edge weights.
  \item \kw{color: categorical(array);} --- discrete classes (coloring, components).
  \item \kw{color: continuous(array);} --- gradient from scores (PageRank, degree).
  \item \kw{size: continuous(array);} or \kw{size: arrayName;} --- vertex radius.
\end{itemize}

Algorithm arrays (same length as $|V|$) are passed by name; the drawer maps
values to visual channels.

\subsection{Full drawing example}

\begin{lstlisting}
graph G {
  directed: true;
  edges: file "draw_weighted_directed.txt";
  TRUE
};

int classes[5] = [1, 2, 1, 3, 2];
real scores[n];

for each vertex v in G {
  scores[v] = v / 4.0;
}

draw G to "auto.svg";

draw G to "hierarchical.pdf" {
  layout: "hierarchical";
  vertices {
    labels: true;
    color: categorical(classes);
    size: continuous(sizes);
  }
  edges { labels: true; }
};
\end{lstlisting}

%==============================================================================
\section{Drawing motif patterns (\texttt{drawmotifs})}
%==============================================================================

Highlight one motif pattern on a host graph:

\begin{lstlisting}
draw motifs of G to "ffl.png" {
  layout: "hierarchical";
  vertices { labels: true; }
  edges { labels: true; }
  a->b;
  a->c;
  b->c;
};
\end{lstlisting}

\ann{%
  Grammar: \texttt{draw ('motif' | 'motifs') 'of' graphID 'to' STRING '\{'
  drawMotifOption* '\}' ';'}.\\
  Motif edge lines inside the block use the same \sym{->} / \sym{-|} syntax as
  comprehension. All matches of the pattern are drawn on \texttt{G}.
}

Signed example:

\begin{lstlisting}
draw motifs of G to "signed_motif.png" {
  layout: "hierarchical";
  a->b;
  a-|c;
  b->c;
};
\end{lstlisting}

%==============================================================================
\section{Complete walk-through: declare, filter, draw, count}
%==============================================================================

The following program ties together declaration, Milo-style motif
comprehension, drawing, and printing (from \texttt{milo\_motif\_examples.graph}):

\begin{lstlisting}
graph G {
  directed: true;
  nodes: 0, 1, 2, 3, 4, 5, 6, 7, 8, 9, 10, 11, 12, 13;
  edges:
    0->1, 0->2, 1->2,
    3->4, 4->5,
    6->8, 6->9, 7->8, 7->9,
    10->11, 10->12, 11->13, 12->13;
};

graph FeedForward = [G where motif {
  a->b; a->c; b->c;
}];

draw G to "original.png" {
  layout: "hierarchical";
  vertices { labels: true; }
};

draw FeedForward to "ffl_subgraph.png" {
  layout: "hierarchical";
  vertices { labels: true; }
};

print numVertices(FeedForward);
print numEdges(FeedForward);
\end{lstlisting}

\ann{%
  \textbf{Step 1} --- \kw{graph G} declares the host network.\\
  \textbf{Step 2} --- \kw{graph FeedForward = [G where motif \{...\}]} extracts
  the union of all feed-forward loop edges.\\
  \textbf{Step 3} --- \kw{draw G} and \kw{draw FeedForward} emit PNG files.\\
  \textbf{Step 4} --- \kw{print numVertices/numEdges} report subgraph size.
}

%==============================================================================
\section{Quick reference: grammar rules to syntax}
%==============================================================================

\begin{longtable}{@{}p{4.2cm}p{9.8cm}@{}}
\toprule
\texttt{Base.g4} rule & Graph-related surface forms \\
\midrule
\endhead
\texttt{graphDef} & \texttt{graph ID \{ properties; nodes; edges; \}} \\
\texttt{graphProperty} & \texttt{directed: true|false;} \\
\texttt{nodes} / \texttt{edges} & \texttt{nodes: 0,1,...;} \quad \texttt{edges: u->v,...;} or \texttt{file "path";} \\
\texttt{graphComprehension} & \texttt{graph H = [G where condition];} \\
\texttt{motifMatchesDecl} & \texttt{motifs M = [G where motif \{...\}];} \\
\texttt{graphListDecl} & \texttt{graphs S = [G where motif \{...\}];} \\
\texttt{graphCondition} & \texttt{degree > k}, \texttt{connected with v}, \texttt{cycle}, \texttt{motif \{...\}} \\
\texttt{foreachStatement} & \texttt{for each vertex|edge|neighbor|graph|motif ... in ...} \\
\texttt{nodeEdgeOperation} & \texttt{add ... to G;} \quad \texttt{remove ... from G;} \\
\texttt{queryStatement} & \texttt{query id : "name" of G;} \\
\texttt{showgraph} & \texttt{show G;} \\
\texttt{drawgraph} & \texttt{draw G to "file.ext" \{ options \};} \\
\texttt{drawmotifs} & \texttt{draw motifs of G to "file" \{ motif; options \};} \\
\texttt{printgraph} & \texttt{print nodes|edges|graph of G;} \\
\texttt{NodeCheck} / \texttt{EdgeCheck} & \texttt{v in G}, \texttt{u->v in G} \\
\bottomrule
\end{longtable}

\vspace{1em}
\noindent
\textbf{Source files for examples:}
\texttt{p2GraphEasy/milo\_motif\_examples.graph},
\texttt{p2GraphEasy/graph\_list\_motif\_test.graph},
\texttt{p2GraphEasy/draw\_visualization\_test.graph},
\texttt{p2GraphEasy/signed\_motif\_test.graph},
\texttt{p2GraphEasy/incoherent\_ffl\_demo.graph},
\texttt{p2GraphEasy/large\_validation/bfs\_parent.graph}.
Formal grammar: \texttt{Base.g4}; prose manual: \texttt{docs/language.md}.

\end{document}
